\section{Discussion}

\subsection{Value of Portfolio-Level Assessment}

The cross-portfolio panel identified insights invisible at the individual paper level:
\begin{itemize}
    \item \textbf{Convergent architecture} across independently-developed modules
    \item \textbf{Mutual validation} that addresses each module's primary weakness
    \item \textbf{Shared gaps} (learning loop, human dimensions) that suggest new research directions
\end{itemize}

These findings demonstrate that portfolio-level review provides qualitatively different---not just aggregated---assessment compared to individual paper review.

\subsection{Structured Disagreement as Signal}

The three reviewer blocs (Systems, Agents, Human-Centered) with near-orthogonal between-bloc correlations provide richer information than consensus. A paper ranked highly by all blocs is genuinely strong across dimensions. A paper ranked highly by one bloc but low by others has a clear audience but limited breadth.

\subsection{Limitations}

\textbf{Panel size}. Seven reviewers is sufficient for structured disagreement analysis but may not capture all relevant perspectives. Larger panels would improve coverage at the cost of coordination.

\textbf{Single program}. The methodology has been applied to one research program (14 papers, 2 modules). Validation across diverse programs is needed.

\textbf{AI-simulated reviewers}. Panel assessments are generated by LLM personas, not real experts. The structured disagreement patterns may reflect prompt engineering rather than genuine disciplinary differences.
